%% ch03 --- Sprzężenie zwrotne: dlaczego nic nie rośnie wiecznie
%%
%% Napisany we wrześniu 2026 według wzoru rozdziału 1. Każda liczba, której
%% czytelnik nie policzy w pamięci, pochodzi z code/measure/ch03.py, który
%% importuje TE SAME funkcje, które drukują listingi (code/ch03/), więc strona
%% i listing nie mogą się rozejść. Rachunki w podrozdziale o dwóch
%% pchnięciach czytelnik wykonuje sam i są wpisane wprost.
%%
%% CZEGO TEN ROZDZIAŁ NIE MOŻE POWIEDZIEĆ. Podkręca tempo wzrostu tylko do
%% pierwszego cyklu o okresie dwa i pokazuje jedynie, że USTALANIE SIĘ USTAJE.
%% Gdzie dokładnie punkt stały traci stabilność, diagram pajęczynowy na
%% paraboli, cykl o okresie cztery i nieregularne zachowanie dalej należą do
%% rozdziału 5; diagram bifurkacyjny do rozdziału 8; r = 4 nie pada wcale.
%% Nachylenie jest MIERZONE, nigdy nie podane wzorem.
%%
%% Bliźniak: chapters/en/ch03-feedback.tex. Podrozdział za podrozdziałem,
%% makro za makrem, liczba za liczbą; tools/parity.py je porównuje.

\chapter{Sprzężenie zwrotne: dlaczego nic nie rośnie wiecznie}
\label{ch:feedback}
\index{sprzężenie zwrotne}\index{wzrost!wykładniczy}

\noindent
W ciepłej szalce z mnóstwem pożywienia bakteria \emph{E.~coli} dzieli się na
dwie mniej więcej co dwadzieścia minut. Zacznij od jednej; każda nowa
komórka robi to samo, a potem każda z jej potomnych. Gdyby nic nie stanęło na
przeszkodzie, po jakim czasie kolonia ważyłaby więcej niż Ziemia?

Nic takiego się nie dzieje, a ten rozdział mówi, dlaczego. Coś zawsze stawia
opór: kończy się pożywienie, zapełnia się przestrzeń, gromadzą się odpady.
Gdy wielkość czegoś zmienia tempo, w jakim to rośnie, mamy do czynienia ze
sprzężeniem zwrotnym, a sprzężenie zwrotne wygina proste reguły z
rozdziału~\ref{ch:iteration} w krzywe. Krzywą trudniej przewidzieć niż
prostą, a reszta tej książki opowiada o tym, o ile trudniej. Tutaj spotkasz
pierwszą niespodziankę: populacja, której kończy się miejsce, nie zawsze
łagodnie się zatrzymuje.

\begin{outcomes}
\outcome{odróżnić dodatnie sprzężenie zwrotne od ujemnego w układzie
  opisanym słowami}
\outcome{zamienić zdanie \enquote{wzrost zwalnia, gdy brakuje miejsca} w
  regułę i ją uruchomić}
\outcome{sprawdzić, czy skutki dwóch pchnięć danej reguły się sumują}
\outcome{znaleźć punkty stałe reguły krzywoliniowej i zmierzyć obliczeniowo
  nachylenie w każdym z nich}
\outcome{przewidzieć na podstawie tego nachylenia, co zrobi populacja w
  pobliżu swojego poziomu ustalonego}
\end{outcomes}

\section{Podwajanie i dlaczego nie może trwać}
\label{sec:ch03-doubling}
\index{podwajanie}\index{bakterie}

\noindent
Podwajanie to najprostszy wzrost, jaki istnieje: na każdym kroku każda
sztuka tego, co masz, zamienia się w dwie. Rozdział~\ref{ch:iteration} nazwał
taki wzrost, proporcjonalny do tego, co już jest, wykładniczym.

\begin{think}
Wróćmy do bakterii, która podwaja się co dwadzieścia minut. Mniej więcej po
jakim czasie kolonia będzie ważyć więcej niż Ziemia: po stu latach, po roku,
po miesiącu czy po dwóch dniach?
\end{think}
\reveal{Po niecałych dwóch dniach: po \val{ch03.earth.doublings}
podwojeniach, czyli po \val{ch03.earth.hours} godziny. Jedna komórka waży
około jednej bilionowej grama.}
Większość ludzi zgaduje rok albo dłużej. Listing liczy na liczbach
całkowitych, więc po drodze nic się nie zaokrągla:

\pyregion{code/ch03/doubling.py}{grow}{Liczenie podwojeń, aż kolonia przeważy Ziemię}{lst:ch03-grow}

\noindent
Wiersz \code{cells = 2 * cells} to cała reguła wzrostu. Reszta programu
wypisuje stan kolonii co osiem godzin:

\pyregion{code/ch03/doubling.py}{table}{Kolonia co osiem godzin}{lst:ch03-table}

\transcript{ch03-doubling}

\noindent
W skrócie, którym Python zapisuje bardzo duże i bardzo małe liczby,
\code{4.7e+09} to \num{4.7} z przecinkiem przesuniętym o dziewięć miejsc w
prawo (Python stawia kropkę tam, gdzie po polsku piszemy przecinek), a
\code{1.0e-12} to przecinek przesunięty o dwanaście miejsc w lewo. Po jednej
dobie kolonia ważyłaby około \val{ch03.day.tonnes} ton; niecałą dobę później
więcej niż planeta.

Podwajanie myli oko, bo każdy krok dodaje więcej niż wszystkie wcześniejsze
razem. Połóż jedno ziarnko ryżu na pierwszym polu szachownicy, dwa na drugim,
cztery na trzecim i tak dalej: pierwsze cztery pola mieszczą
$1 + 2 + 4 + 8 = 15$ ziarenek, a piąte samo $16$, i na końcu ostatnie pole
mieści więcej niż pozostałe sześćdziesiąt trzy razem.

Nic nie rośnie tak długo. W kolbie z bulionem bakterie zatrzymują się na
kilku miliardach komórek w każdym mililitrze, gdy kończy się pożywienie i
zbierają się odpady. Wzrost, który czymś się żywi, natrafia na granicę tego
czegoś.

\section{Sprzężenie zwrotne}
\label{sec:ch03-feedback}
\index{sprzężenie zwrotne!dodatnie}\index{sprzężenie zwrotne!ujemne}

\noindent
Spójrz jeszcze raz na regułę podwajania. Im więcej jest bakterii, tym więcej
podziałów zachodzi w ciągu następnych dwudziestu minut: wielkość kolonii
ustala jej własne tempo wzrostu. Gdy wielkość czegoś zmienia jego własne
tempo zmian, mówimy o \emph{sprzężeniu zwrotnym}: skutek wraca i działa na
to, co go wywołało.

Sprzężenie zwrotne bywa dwojakie. \emph{Dodatnie sprzężenie zwrotne} pcha
dalej w kierunku, w którym sprawy już zmierzają. Zbliż mikrofon do głośnika,
do którego jest podłączony, a cichy dźwięk zacznie krążyć w pętli, za każdym
razem głośniejszy, aż wycie osiągnie największą głośność, na jaką stać
wzmacniacz. \emph{Ujemne sprzężenie zwrotne} stawia opór. Termostat wyłącza
ogrzewanie, gdy w pokoju jest za ciepło, i włącza je, gdy jest za zimno.
Zatłoczenie to także ujemne sprzężenie zwrotne: im więcej zwierząt, tym mniej
pożywienia dla każdego i tym mniej młodych w następnym pokoleniu. Drapieżniki
i ofiary tworzą podwójną pętlę: więcej królików wykarmi więcej lisów, więcej
lisów zje więcej królików, a mniej królików zagłodzi lisy.

\mermaidfig{ch03-crowding-loop}{Zatłoczenie jako ujemne sprzężenie zwrotne.
Im więcej jest zwierząt, tym mniej młodych zostawia każde z nich, więc pętla
ściąga populację z powrotem.}{zatłoczenie jako ujemna pętla}

\begin{think}
Plotka głosi, że bank może upaść, i część klientów wypłaca pieniądze. Czy to
popycha bank ku upadkowi, czy od niego oddala \dash{} a więc czy to
sprzężenie zwrotne dodatnie, czy ujemne?
\end{think}
\reveal{Dodatnie. Każda wypłata osłabia bank, a to sprawia, że następny
klient chętniej też wypłaci pieniądze. Run na bank sam się napędza.}
Run na bank Northern Rock we wrześniu 2007 roku, pierwszy na brytyjski bank
od ponad stu lat, skończył się dopiero wtedy, gdy rząd zagwarantował
depozyty: ujemne sprzężenie zwrotne, dodane z zewnątrz. Dodatnie sprzężenie
kończy się właśnie tak albo na jakiejś granicy, a od granicy zaczyna się
ciekawa matematyka.

\section{Wzrost zwalnia, gdy brakuje miejsca}
\label{sec:ch03-logistic}
\index{wzrost logistyczny}\index{Verhulst, Pierre-François}\index{krzywa S}

\noindent
Oto najprostszy sposób, by wpisać zatłoczenie w regułę. Mierz populację nie
w zwierzętach, lecz jako ułamek tego, ile to miejsce mogłoby kiedykolwiek
pomieścić. Nazwij go $x$: $x = 0$ to miejsce puste, a $x = 1$ to takie, w
którym zajęty jest każdy okruch pożywienia i przestrzeni. Wtedy $1 - x$ to
część miejsca, która jest jeszcze wolna.

Gdy miejsca jest pod dostatkiem, przyjmij, że każde zwierzę zostawia $r$
młodych, które dożyją rozrodu, więc następne pokolenie jest $r$ razy
liczniejsze od obecnego: $x \to r x$. Dla $r = 2$ to reguła podwajania.
Teraz niech zatłoczenie ją spowolni: pomnóż przez część miejsca, która jest
jeszcze wolna:
\[ x \;\to\; r\,x\,(1 - x). \]
Słowami: następne pokolenie to tyle, ile obecne wydałoby na świat, mając
miejsca pod dostatkiem, pomniejszone w proporcji do miejsca, które zostało.
Dopóki $x$ jest maleńkie, $1 - x$ jest bliskie $1$ i reguła jest prawie
regułą podwajania. Gdy $x$ rośnie, rośnie z nim hamulec, a przy $x = 1$
następnego pokolenia nie ma wcale. Belgijski matematyk Pierre-François
Verhulst zaproponował prawo wzrostu, które zwalnia w ten sposób, w 1838
roku; powyższa reguła to jego pomysł rozpisany na kolejne pokolenia.

\pyregion{code/ch03/crowding.py}{rule}{Reguła zatłoczenia, pokolenie po pokoleniu}{lst:ch03-rule}

\noindent
Uruchom ją obok zwykłego podwajania, obie od jednej tysięcznej miejsca, przy
$r = 2$:

\pyregion{code/ch03/crowding.py}{run}{Podwajanie oraz podwajanie, gdy brakuje miejsca}{lst:ch03-run}

\transcript{ch03-crowding}

\noindent
Przez sześć pokoleń kolumn prawie nie da się odróżnić; miejsce zapełnione w
kilku procentach nie wydaje się zatłoczone. Potem się rozchodzą. Kolumna
podwajania przekracza całe miejsce w pokoleniu \val{ch03.exp.full}, co jest
niemożliwe, i rośnie dalej, jakby nigdy nic. Kolumna zatłoczenia zwalnia,
ustala się na \num{0.5} i tam zostaje.

\begin{think}
Przyjrzyj się kolumnie zatłoczenia. W których pokoleniach populacja
przybiera najwięcej: gdy jest mała, gdy jest mniej więcej w połowie drogi do
miejsca, w którym się zatrzyma, czy gdy jest już prawie u celu?
\end{think}
\reveal{Mniej więcej w połowie drogi. Największy skok jest z pokolenia
\val{ch03.fast.from} do pokolenia \val{ch03.fast.to}: populacja przybiera
wtedy \val{ch03.fast.gain} miejsca i mija \val{ch03.fast.level}, połowę
poziomu \num{0.5}, na którym się ustala.}
Mała populacja ma niewiele zwierząt, które mogłyby się rozmnażać; prawie
pełna nie ma miejsca, żeby się rozmnażać. Sprawdzenie każdego poziomu od $0$
do $1$ co jedną stutysięczną potwierdza, że przyrost w jednym pokoleniu jest
największy dokładnie w połowie poziomu ustalonego. Na wykresie kolumna
zatłoczenia układa się w krzywą w kształcie litery S, krzywą wzrostu, który
natrafia na granicę.

\plotfig{ch03-s-curve}{Podwajanie oraz podwajanie, gdy brakuje miejsca, oba
od jednej tysięcznej miejsca. Zgadzają się, dopóki miejsce jest prawie
puste; potem jedno ucieka poza wykres, a drugie wygina się w S.}{wzrost
wykładniczy a logistyczny}

\begin{worldbox}
Nauka o rybołówstwie używa krzywych wzrostu tego kształtu od połowy XX
wieku. Przełowiona populacja ryb odbudowuje się najszybciej, gdy ma mniej
więcej połowę liczebności sprzed połowów, więc największy połów, jaki można
brać rok po roku bez uszczuplania populacji, \emph{maksymalny podtrzymywalny
połów}, to ten, który utrzymuje ją blisko tej połowy. Ta idea jest zapisana
w Konwencji Narodów Zjednoczonych o prawie morza z 1982 roku. To wciąż
tylko model, a prawdziwe populacje ryb upadały przy limitach, które na
papierze wyglądały bezpiecznie.
\end{worldbox}

\section{Dwa pchnięcia, które się nie sumują}
\label{sec:ch03-superposition}
\index{superpozycja}\index{nieliniowość}

\noindent
Każda reguła z rozdziału~\ref{ch:iteration} była prostą: narysuj następną
wartość w zależności od obecnej, a dostaniesz linię prostą. Reguły
prostoliniowe mają własność tak naturalną, że pewnie przypisujesz ją każdej
regule: skutki pchnięć się sumują. Weź regułę podwajania i populację
\num{0.2}. Dorzuć dodatkowe \num{0.1}, a następne pokolenie wzrośnie o
\num{0.2}; dorzuć \num{0.2}, a wzrośnie o \num{0.4}. Dwa pchnięcia naraz
robią to, co każde z osobna, zsumowane. Nazywa się to \emph{superpozycją} i
to ona sprawia, że reguły prostoliniowe są łatwe: problem można rozłożyć na
części, rozwiązać każdą osobno i dodać odpowiedzi.

Teraz reguła zatłoczenia, z $r = 2$ i tą samą populacją \num{0.2}. Bez
pchnięcia następne pokolenie to $2 \times \num{0.2} \times \num{0.8} =
\num{0.32}$. Dorzuć \num{0.1}, a będzie to $2 \times \num{0.3} \times
\num{0.7} = \num{0.42}$: pchnięcie podniosło je o \num{0.1}.

\begin{think}
Teraz dorzuć \num{0.2}, dwa takie pchnięcia naraz. O ile wzrośnie następne
pokolenie? Najpierw zgadnij, potem policz $2 \times \num{0.4} \times
\num{0.6}$.
\end{think}
\reveal{O \num{0.16}, nie o \num{0.2}: następne pokolenie to \num{0.48}.
Drugie pchnięcie było warte tylko \num{0.06}, bo trafiło na bardziej
zatłoczone miejsce niż pierwsze.}
Wyżej robi się jeszcze dziwniej. Przy populacji \num{0.6} następne pokolenie
to $2 \times \num{0.6} \times \num{0.4} = \num{0.48}$; dorzuć \num{0.1}, a
będzie to $2 \times \num{0.7} \times \num{0.3} = \num{0.42}$. Dodanie zwierząt
sprawia teraz, że następne pokolenie jest \emph{mniejsze}. To samo pchnięcie
pomaga na dole i szkodzi na górze: skutek zmiany zależy od tego, gdzie już
jesteś. To właśnie znaczy, że reguła jest \emph{nieliniowa}, i to właśnie
sprzężenie zwrotne robi z regułą.

\begin{trapbox}
\textbf{\enquote{Skutek dwóch pchnięć to suma ich skutków.}} Tylko dla
reguły prostoliniowej. Na krzywej skutek pchnięcia zależy od tego, gdzie
układ już się znajduje, i od wszystkiego, co jeszcze go pcha, więc dwie
przyczyny, z których każda robi niewiele, mogą razem zrobić znacznie więcej,
znacznie mniej albo coś przeciwnego. Każda reguła ze sprzężeniem zwrotnym
jest krzywa.
\end{trapbox}

\section{Gdzie populacja się ustala, i nachylenie}
\label{sec:ch03-slope}
\index{punkt stały}\index{nachylenie}\index{stabilność}

\noindent
Populacja w zatłoczeniu ustaliła się na \num{0.5}. Dlaczego właśnie tam?
Populacja stoi w miejscu, gdy następne pokolenie jest tej samej wielkości co
obecne. Jednym takim poziomem jest puste miejsce: z niczego nic się nie
urodzi. Drugim jest poziom, na którym każde zwierzę zostawia średnio
dokładnie jedno młode na swoje miejsce. Gdy wolna jest część $1 - x$
miejsca, każde zwierzę zostawia $r(1 - x)$ młodych, a to jest jeden, gdy
$1 - x$ wynosi $1/r$, czyli przy
\[ x = 1 - \frac{1}{r}. \]
Dla $r = 2$ to $1 - \tfrac{1}{2} = \tfrac{1}{2}$, tam, gdzie ustaliła się
kolumna. Poziom, którego reguła nie zmienia, to \emph{punkt stały}, jak w
rozdziale~\ref{ch:iteration}.

\begin{think}
Gdzie ustali się populacja przy $r = \num{1.5}$?
\end{think}
\reveal{Na jednej trzeciej: $1/\num{1.5}$ to dwie trzecie, a
$1 - \tfrac{2}{3} = \tfrac{1}{3}$.}
Punkt stały to tylko połowa odpowiedzi; druga połowa to pytanie, czy
populacja w jego pobliżu jest do niego przyciągana, czy od niego odpychana.
Dla reguły prostoliniowej $x \to a x + b$ rozstrzygał to mnożnik $a$:
szturchnięcie od punktu stałego jest na każdym kroku mnożone przez $a$, więc
maleje, gdy $a$ leży między $-1$ a $1$, a w przeciwnym razie rośnie. Krzywa
nie ma jednego mnożnika; jej stromość zmienia się od miejsca do miejsca. Ale
przybliż ją, a w pobliżu dowolnego punktu wygląda jak prosta \dash{} a prosta
ma mnożnik. To jej stromość, czyli \emph{nachylenie}.

\plotfig{ch03-zoom}{Reguła zatłoczenia przy $r = \num{1.5}$, przybliżona
wokół punktu stałego w jednej trzeciej. Z bliska krzywej nie da się odróżnić
od prostej, a stromość tej prostej rozstrzyga, co stanie się z małym
szturchnięciem.}{przybliżanie, aż krzywa wygląda na prostą}

\noindent
Zmierz więc nachylenie: szturchnij populację o jedną milionową w górę i o
jedną milionową w dół i zobacz, jak daleko od siebie wylądują następne
pokolenia.

\pyregion{code/ch03/settle.py}{slope}{O ile jedno pokolenie mnoży maleńkie szturchnięcie}{lst:ch03-slope}

\pyregion{code/ch03/settle.py}{table}{Nachylenie w obu punktach stałych, dla dwóch wartości tempa wzrostu}{lst:ch03-slopes}

\transcript{ch03-settle}

\noindent
W punkcie stałym równym jednej trzeciej nachylenie wynosi
\val{ch03.slope.third}: małe szturchnięcie co pokolenie maleje o połowę, więc
populacja jest przyciągana z powrotem, a punkt stały jest \emph{stabilny}. W
pustym miejscu nachylenie wynosi \val{ch03.slope.zero}: kilka zwierząt w
pustym miejscu się mnoży, czyli rośnie, i ten punkt stały jest
\emph{niestabilny}. Przy $r = 2$ nachylenie w jednej drugiej wynosi
\val{ch03.slope.half}. Ten punkt stały leży na płaskim szczycie garbu, gdzie
szturchnięcie znika prawie całkowicie w jednym pokoleniu, i dlatego kolumna
zatłoczenia tak szybko wskoczyła na \num{0.5}.

Na tej idei opiera się reszta książki: \textbf{na krzywej nachylenie w danym
punkcie jest mnożnikiem maleńkiego szturchnięcia w tym miejscu}, a znajduje
się je, szturchając i patrząc.

\begin{rigourbox}
Przyjęte tu na wiarę: że nachylenie rozstrzyga dla krzywej to, co mnożnik
rozstrzygał dla prostej. Jest tak dla szturchnięć tak małych, że krzywa
wygląda na prostą, co przybliżenie pokazuje, ale czego nie dowodzi; dowód
jest w każdym pierwszym kursie układów dynamicznych. Pomiar szacuje to, co
rachunek różniczkowy nazywa \emph{pochodną}, i tutaj jest dokładny co do
każdej wydrukowanej cyfry.
\end{rigourbox}

\begin{lab}{l03_01_slope}{Zmierz nachylenie i oceń punkt stały}
Napisz \code{slope}, które mierzy nachylenie dowolnej reguły, jaką mu
podasz; \code{fixed\_points}, które podaje oba punkty stałe reguły
zatłoczenia; oraz \code{is\_stable}, które na podstawie nachylenia mówi, czy
punkt stały trzyma. Testy obejmują przycisk cosinusa na kalkulatorze,
naciskany raz za razem.
\end{lab}

\begin{projectbox}
Wspólna biblioteka książki, \pkg{chaoslab}, ma ten pomiar dla dowolnej
reguły jako \api{slope}, a późniejsze rozdziały z niego korzystają.
\end{projectbox}

\section{Przestrzelenie}
\label{sec:ch03-overshoot}
\index{przestrzelenie}\index{oscylacja!tłumiona}\index{opóźnienie}

\noindent
Podkręć tempo wzrostu do $r = \num{2.8}$. Poziom ustalony to
$1 - 1/\num{2.8}$, około \val{ch03.over.level}.

\begin{think}
Zacznij od małej populacji, jak wcześniej. Czy podpełznie do
\val{ch03.over.level} i się zatrzyma, tak jak populacje przy
$r = \num{1.5}$ i $r = 2$?
\end{think}
\reveal{Nie. Wznosi się ponad ten poziom, do \val{ch03.over.peak} w
pokoleniu \val{ch03.over.peakgen}, spada poniżej, do \val{ch03.over.fall},
i waha się w tę i z powrotem, a każde wahnięcie jest mniejsze od
poprzedniego, zanim się ustali.}

\begin{trapbox}
\textbf{\enquote{Granica wzrostu oznacza, że wzrost łagodnie się
zatrzymuje.}} Może tak być i przy łagodnym tempie wzrostu tak jest. Przy
silnym populacja przestrzeliwuje poziom, który otoczenie jest w stanie
utrzymać, zapada się poniżej niego i oscyluje. Granica wciąż istnieje, ale
dochodzenie do niej to seria wahań, a po dalszym podkręceniu wahania wcale
nie muszą wygasać.
\end{trapbox}

\noindent
Nachylenie mówi dlaczego. W poziomie ustalonym przy $r = \num{2.8}$ wynosi
ono \val{ch03.slope.over}. Co do wielkości jest mniejsze od jedności, więc
szturchnięcie maleje i poziom jest stabilny. Ale jest ujemne: populacja
trochę powyżej poziomu zostaje posłana trochę poniżej niego, potem znów
trochę powyżej. Ujemne nachylenie to przestrzelenie wbudowane w regułę, a
ustalanie się przez coraz mniejsze wahania nazywa się \emph{oscylacją
tłumioną}.

Przestrzelenie bierze się z opóźnienia. Reguła wylicza zmianę na całe
pokolenie z tego, jak zatłoczone jest miejsce teraz, i patrzy znowu dopiero
pokolenie później. Przy \val{ch03.over.before}, poniżej swojego poziomu,
populacja widzi mnóstwo miejsca i wydaje na świat liczne następne pokolenie;
zanim te młode się pojawią, miejsca już nie ma, a populacja jest ponad
poziomem. Znasz to z obcego prysznica: przekręcasz kurek w stronę gorącej
wody, nic się nie dzieje, więc przekręcasz dalej \dash{} a potem cię parzy.

Żeby sprawdzić to wyjaśnienie, usuń opóźnienie. Zachowaj to samo prawo
zmiany, ale niech populacja sprawdza swoje zatłoczenie dziesięć razy na
pokolenie: podziel zmianę w każdym pokoleniu na dziesięć rat, każdą
wyliczaną z populacji takiej, jaka jest po poprzedniej racie.

\pyregion{code/ch03/overshoot.py}{instalments}{To samo prawo zmiany, spłacane w dziesięciu ratach}{lst:ch03-instalments}

\pyregion{code/ch03/overshoot.py}{compare}{Raz na pokolenie wobec dziesięciu razy na pokolenie}{lst:ch03-compare}

\transcript{ch03-overshoot}

\noindent
Obie zmierzają do tego samego poziomu, bo prawo zmiany jest to samo i znika
w tym samym miejscu, ale przestrzeliwuje tylko populacja liczona raz na
pokolenie. Ta spłacana w ratach na początku rośnie nawet szybciej, bo jej
młode zaczynają się rozmnażać jeszcze w tym samym pokoleniu, a mimo to
wsuwa się w swój poziom, nie mijając go. Przestrzelenie powoduje nie
szybkość wzrostu, lecz to, jak późno działają hamulce. Własne prawo
Verhulsta to to, czym stają się raty, gdy robi się je coraz mniejsze, i ono
nigdy nie przestrzeliwuje.

\plotfig{ch03-overshoot}{To samo prawo wzrostu przy $r = \num{2.8}$. Liczone
raz na pokolenie przestrzeliwuje swój poziom i się waha; liczone dziesięć
razy na pokolenie wsuwa się w ten sam poziom.}{przestrzelenie i to samo
prawo bez opóźnienia}

\needspace{10\baselineskip}
\begin{lab}{l03_02_instalments}{Późne hamulce i wczesne hamulce}
Napisz regułę zatłoczenia, tę samą regułę spłacaną w \code{k} ratach,
funkcję, która uruchamia każdą z nich na wiele pokoleń, oraz taką, która
mówi, czy populacja kiedykolwiek mija swój poziom ustalony. Potem wypróbuj
własne wartości tempa wzrostu.
\end{lab}

\section{Kiedy populacja przestaje się ustalać}
\label{sec:ch03-cycle}
\index{cykl}

\noindent
Podkręć tempo wzrostu jeszcze raz, do $r = \num{3.1}$. Poziom ustalony wciąż
istnieje: $1 - 1/\num{3.1}$ to około \val{ch03.two.level}, a populacja
umieszczona dokładnie na nim by na nim została. Ale nachylenie w tym
miejscu, zmierzone jak wcześniej, wynosi \val{ch03.slope.two}.

\begin{think}
Co nachylenie \val{ch03.slope.two} robi z małym szturchnięciem od poziomu
ustalonego, pokolenie po pokoleniu?
\end{think}
\reveal{W każdym pokoleniu przerzuca szturchnięcie na drugą stronę i
powiększa je o jedną dziesiątą. Szturchnięcie rośnie, więc poziom ustalony
nie jest już w stanie utrzymać populacji.}

\pyregion{code/ch03/two_cycle.py}{nudge}{Szturchnięcie o jedną tysięczną, pokolenie po pokoleniu}{lst:ch03-nudge}

\pyregion{code/ch03/two_cycle.py}{run}{Gdzie populacja ląduje}{lst:ch03-end}

\transcript{ch03-two-cycle}

\noindent
Szturchnięcie rośnie, co pokolenie zmieniając stronę, a populacja ląduje
daleko od poziomu ustalonego. Waha się między dwiema wartościami,
\val{ch03.two.low} i \val{ch03.two.high}, w jednym pokoleniu nisko, w
następnym wysoko, bez końca. Nie pojawiło się nic losowego; reguła to wciąż
ta sama jedna linijka. Ustalanie się ustało, a jego miejsce zajął cykl, który
powtarza się co dwa pokolenia. Populacja spłacana w ratach wciąż ustala się
na \val{ch03.two.level}: usuń opóźnienie, a wahania znikną razem z nim.

Gdzieś między $r = \num{2.8}$ a $r = \num{3.1}$ poziom ustalony stracił
swoją moc przyciągania. Gdzie dokładnie i co robi reguła, gdy tempo wzrostu
podkręca się jeszcze dalej, to temat rozdziału~\ref{ch:logistic}. Reguła ma
jedną linijkę, dziecko policzyłoby ją ołówkiem, a jeszcze nie skończyła cię
zaskakiwać.

\begin{summarybox}
\begin{itemize}
\item \textbf{Wzrost wykładniczy}, proporcjonalny do tego, co już jest, na
  każdym kroku dodaje więcej niż wszystkie wcześniejsze kroki, i nic
  rzeczywistego nie rośnie tak długo.
\item \textbf{Sprzężenie zwrotne} to wielkość, która zmienia własne tempo
  zmian. \textbf{Dodatnie sprzężenie zwrotne} pcha dalej (wycie głośnika, run
  na bank), aż natrafi na granicę; \textbf{ujemne sprzężenie zwrotne} stawia
  opór (termostat, zatłoczenie).
\item Zatłoczenie wpisane w podwajanie daje $x \to r x (1 - x)$, gdzie $x$ to
  zajęta część miejsca: \textbf{krzywą S}, rosnącą najszybciej w połowie
  poziomu ustalonego.
\item Dla reguły prostoliniowej skutki pchnięć się sumują
  (\textbf{superpozycja}); dla krzywej, \textbf{nieliniowej} reguły nie, a
  pchnięcie może pomóc albo zaszkodzić, zależnie od tego, gdzie jest układ.
\item Reguła ustala się tam, gdzie każde zwierzę jest dokładnie zastępowane,
  w \textbf{punkcie stałym} $1 - 1/r$. \textbf{Nachylenie} w nim, zmierzone
  szturchaniem, mnoży małe szturchnięcie; wielkość mniejsza od jedności
  oznacza stabilność.
\item Ujemne nachylenie oznacza \textbf{przestrzelenie} i \textbf{oscylację
  tłumioną}. Przyczyną jest opóźnienie: to samo prawo w małych ratach nie
  przestrzeliwuje.
\item Po dalszym podkręceniu wielkość nachylenia przekracza jeden, a
  populacja bez końca waha się między dwiema wartościami \dash{} to wstęp do
  rozdziału~\ref{ch:logistic}.
\end{itemize}
\end{summarybox}

\canyou

\begin{checkpoint}
\item Ile ziarenek leży na jedenastym polu szachownicy, a ile na pierwszych
  dziesięciu polach razem?
  \answerto{$1024$ na jedenastym i $1023$ na pierwszych dziesięciu: każde pole
  mieści o jedno ziarenko więcej niż wszystkie pola przed nim.}
\item Dodatnie czy ujemne sprzężenie zwrotne: termostat; oklaski, które stają
  się głośniejsze, gdy dołączają kolejne osoby; cena, która rośnie, gdy
  towaru brakuje, więc mniej ludzi go kupuje?
  \answerto{Ujemne, dodatnie, ujemne.}
\item Co w regule zatłoczenia przy $r = 2$, zaczynając od \num{0.5}, robią z
  następnym pokoleniem pchnięcia o \num{0.1} i o \num{0.2}?
  \answerto{Zmienia się ono z \num{0.5} na \num{0.48} i na \num{0.42}: spada
  o \num{0.02} i o \num{0.08}. Dwa razy większe pchnięcie wyrządza cztery
  razy większą szkodę.}
\item Gdzie ustala się populacja przy $r = \num{2.5}$? A przy
  $r = \num{1.25}$?
  \answerto{Na \num{0.6} i na \num{0.2}: $1 - 1/\num{2.5}$ oraz
  $1 - 1/\num{1.25}$.}
\item Nachylenie w punkcie stałym wynosi $-\num{0.5}$, a populacja startuje
  \num{0.01} nad nim. Gdzie jest po jednym pokoleniu, a gdzie po dwóch?
  \answerto{\num{0.005} pod punktem stałym, potem \num{0.0025} nad nim.
  Ustala się, przestrzeliwując, a każde wahnięcie jest o połowę mniejsze od
  poprzedniego.}
\item Dlaczego populacja, która sprawdza swoje zatłoczenie dziesięć razy na
  pokolenie, nie przestrzeliwuje, skoro ta, która sprawdza je raz na
  pokolenie, przestrzeliwuje?
  \answerto{Reaguje na zatłoczenie w miarę, jak narasta, a nie z opóźnieniem
  całego pokolenia: przestrzelenie bierze się z opóźnienia, nie z szybkości
  wzrostu.}
\end{checkpoint}
